\documentclass[11pt,a4paper]{article}
\usepackage[margin=2.4cm]{geometry}
\usepackage{booktabs,longtable,array}
\usepackage[hidelinks]{hyperref}
\usepackage{xcolor}
\usepackage{enumitem}
\newcommand{\todo}[1]{\textcolor{red}{[TBD: #1]}}
\newcommand{\claim}[1]{\textbf{#1}}
\usepackage{tabularx} % Make sure to add this to your preamble
\usepackage{booktabs}
\title{Pre-registered protocol: a global sensitivity audit of the
scenario-generation pipeline of Li et al.\ (2025), \emph{Nature Climate
Change}\\[4pt]\large Modelling of the Modelling Process (MoMP) series ---
optimal-transport GSA case}
% \author{Samuele Lo Piano$^{1}$ \and Giulia Vannucci$^{2}$ \and
% Arnald Puy$^{3}$ \and Roberta Siciliano$^{2}$ \and Andrea Saltelli$^{4}$
% \\[6pt]\small $^{1}$Gda\'nsk University of Technology \quad
% $^{2}$University of Naples Federico II \quad
% $^{3}$University of Birmingham \quad $^{4}$University of Bergen
% \\\small \todo{confirm author list and order with co-authors}}
\date{Version 0.1 --- \todo{date of Zenodo deposit} --- DOI: \todo{minted at deposit}}

\begin{document}
\maketitle

\begin{abstract}
\noindent We pre-register a global sensitivity analysis (GSA) of the
scenario-generation pipeline of Li, Zhu, McJeon, Byers, Zhou \& Ou (2025),
``Using deep learning to generate key variables in global mitigation
scenarios'', \emph{Nat.\ Clim.\ Change} 15:760--768
(doi:10.1038/s41558-025-02352-8). The object of analysis is not the climate
or the AR6 ensemble but the pipeline itself: the analyst's methodological
choices are the input factors, the generated multivariate ensemble and the
paper's validation claims are the outputs, and sensitivity is measured by
optimal-transport (Wasserstein) displacement of the synthetic ensemble,
estimated from given data. This protocol freezes, before any GSA run: the
operational rules for the paper's claims (\S\ref{sec:claims}), the factor
space (\S\ref{sec:factors}), the experimental design and compute budget
(\S\ref{sec:design}), the estimators and inference procedure
(\S\ref{sec:estimators}), and the decision criteria (\S\ref{sec:criteria}).
The completed replication record (Phase~0, \S\ref{sec:phase0}) is reported
for transparency. At the time of this deposit, no claim rule
(X1a--X5) has been evaluated at any factor setting, including the paper's
own values, no retraining has been performed, and no noise-floor replicate
has been run. The only pipeline stage executed is the deterministic
regeneration check of Phase~0B, which evaluates no claim.
\end{abstract}

\section{Purpose and scope}
Continuing the MoMP programme, this study asks whether the properties Li
et al.\ attribute to their synthetic-scenario framework are properties of
the framework or of the specific, partly undocumented methodological
choices through which it was instantiated. We audit the epistemic
apparatus, not the climate--economy content: factors are analyst choices,
outputs are (i) the synthetic ensemble as a multivariate distribution and
(ii) pass/fail evaluations of the paper's claims under fixed operational
rules.

\section{Audited artefacts}\label{sec:artefacts}
All artefacts are pinned by version and SHA-256 in \texttt{DATA.md} of the
project repository (archived with this deposit). Principal items:
\begin{itemize}[nosep]
\item \textbf{A1} Deep-IAM code, v1.0.0, doi:10.5281/zenodo.15280653 ---
  the version cited by the paper, the audited object.
\item \textbf{A2} Deep-IAM v1.0.1 (2026-04-25). Admitted as documented
  supplement only for its two added files:
  \texttt{Secondary-data-generate.ipynb} (the generator of the deposited
  synthetic data; absent from v1.0.0) and \texttt{Configuration of
  Environments.txt} (environment pins: Python 3.8.20, TensorFlow 2.13.0,
  tsgm 0.0.7, scikit-learn 1.3.2, XGBoost 2.1.1). Full diff archived.
\item \textbf{A3} Synthetic-scenario deposit, doi:10.5281/zenodo.15240553
  (Gen/Real CSVs; validated in Phase~0).
\item \textbf{B1--B2} AR6 Scenarios Database v1.1 (metadata + World
  timeseries), doi:10.5281/zenodo.7197970.
\end{itemize}

\section{Phase 0 --- completed replication record}\label{sec:phase0}
Completed before this deposit; machine-readable reports archived under
\texttt{results/}. Summary:
\begin{enumerate}[nosep]
\item \textbf{Metadata fidelity.} AR6 v1.1 vetted-with-climate population
  reproduces the paper's Extended Data Table~1 exactly: $n=1202$;
  per-category counts match digit for digit (\texttt{phase0\_metadata\_check.json}).
\item \textbf{First undeclared filter reconstructed (D2).} The training
  population (SI~S3.1: 680/308/172, $n=1160$) is reproduced exactly by:
  complete 2020--2100 coverage on the five feature variables
  (\texttt{Primary Energy|Coal/Oil/Gas}, \texttt{Final Energy|Liquids},
  \texttt{Carbon Sequestration|CCS}); Kyoto Gases exempt. This rule appears
  nowhere in the paper or SI.
\item \textbf{Second undeclared filter reconstructed (D9).} The population
  underlying the paper's figures and label-transfer benchmark (Fig.~4a,
  ED~Fig.~2, Table~1: 554/237/110, $n=901$) is reproduced exactly by:
  complete 2020--2100 coverage on the five feature variables and all
  ten electricity fuel-level variables, Kyoto Gases exempt, raw (pre-imputation)
  CS availability required. Across the two silent drops the smallest
  category (C78) loses 43\% of its vetted scenarios (193$\to$110).
\item \textbf{Deposit validated.} After harmonizing a Mt-vs-Gt unit
  inconsistency (deposit in Mt, the paper's Source Data Fig.~3 in Gt), all
  1{,}620 quantile comparisons (6 variables $\times$ 6 series $\times$ 9
  years $\times$ 5 quantiles) pass at 2\% relative tolerance; Gen files
  hold 10{,}000 scenarios per category.
\item \textbf{Version provenance (D7).} v1.0.0 cannot reproduce the
  deposited synthetic data or Fig.~5. The generating notebook appeared in
  v1.0.1 ten months after publication, without changelog.
\item \textbf{Non-reproducibility of training (D1).} The authors state that
  TSGM training is not reproducible even under fixed seeds and distribute
  trained weights, this motivates the noise-floor design of
  \S\ref{sec:design}.
\item \textbf{Phase 0B --- executability and determinism of the generation
  stage (completed).} Evaluating no claim rule, we executed the pipeline's
  one nominally deterministic link: the authors' shipped decoder weights
  (\texttt{vae\_conv5}, 9 steps $\times$ 16 features, latent dim.\ 10) were
  loaded, the \texttt{TSFeatureWiseScaler} refit on the deposited Real CSVs
  (which the D9 reconstruction identifies as the training tensors), the
  latent noise drawn under the documented global seed 3, decoded, and
  compared to the deposited Gen CSVs (environment: Python 3.10.21,
  TensorFlow 2.13.0, tsgm 0.0.7, macOS/arm64; the authors used Windows
  x86-64, Python 3.8). Results
  (\texttt{results/phase0b\_regeneration.json}): (i)~\emph{executability
  confirmed} --- weights load, architectures match, generation runs
  end-to-end from the archived artefacts alone; (ii)~\emph{sample-level
  reproduction fails} under all three latent-stream readings of the
  documentation (TensorFlow global-seed, a second-execution probe, and the
  SI's NumPy reading): median elementwise relative deviation 0.33--0.45,
  i.e.\ entirely different draws; (iii)~\emph{distributional agreement in
  the bulk}: median quantile deviation 1.8\% of the reference spread
  (Q90--Q10), 95th percentile 10.3\%, yet $\approx 4.4\times$ above the
  Monte-Carlo floor obtained from two independent 10{,}000-draw ensembles
  through the same decoder, with the excess concentrated in distribution
  tails and tight-spread cells (e.g.\ carbon sequestration in 2020);
  (iv)~a scaler-bounds diagnostic flags no feature, consistent with the
  Real CSVs being the scaler-fit data. The residual is most plausibly a
  checkpoint mismatch (the shipped weights differing from those that
  generated the deposit) consistent with the repository README's account
  of multiple training sessions with ``comparable'' results. Conclusion:
  the generation stage is executable and distributionally faithful, but
  the deposited outputs cannot be tied to the documented seed or,
  apparently, to the shipped weights; the determinism gap the authors
  acknowledge for training (D1) extends to generation (catalogued as D11).
  This check answers no question posed in
  \S\ref{sec:claims}--\S\ref{sec:criteria}.
\end{enumerate} \\
The full inconsistency catalogue (D1--D10) is archived as
\texttt{DISCREPANCIES.md}.

\section{Claims and operational rules}\label{sec:claims}
Each claim of the target paper is translated into a rule verifiable on the
pipeline's output. Populations are declared explicitly, given D2/D9.
Thresholds marked \todo{} are to be fixed at freeze. No rule will be
altered after deposit except under the deviations policy
(\S\ref{sec:deviations}).

\begin{description}[style=nextline]
\item[\claim{X1a} (distributional match).] ``[The VAE] closely matches the
  original AR6 scenarios in terms of mean and the spread, particularly for
  C1234.'' \emph{Rule:} for each of the six key variables and each
  category, the median over years of the absolute relative deviation
  between synthetic and real quantiles (Q10, Q25, Q50, Q75, Q90; real
  population $=$ the 901-set) is $\le \tau_{1}$, with
  $\tau_{1}=$ \todo{propose 0.15}; ``particularly C1234'' holds if the
  C1234 deviation is the smallest of the three categories.
\item[\claim{X1b} (narrower distributions).] ``Synthetic scenarios have
  narrower distributions for all climate categories.'' \emph{Rule:}
  IQR$_{\mathrm{syn}} <$ IQR$_{\mathrm{real}}$ for all six variables and
  all three categories at 2050 and 2100.
\item[\claim{X2} (correlation replication).] ``The synthetic scenarios also
  retained the correlations between key variables \dots matching closely
  those observed in the original AR6 scenarios.'' \emph{Rule:} (i) the
  maximum absolute elementwise difference between synthetic and real
  correlation matrices (six variables; 2030, 2050, 2100) is
  $\le \tau_{2}=$ \todo{propose 0.20}; (ii) the paper's reported temporal
  ordering (coal--Kyoto correlation stronger in 2030 than 2100) holds in
  the synthetic data.
\item[\claim{X3} (internal consistency).] ``Among the 30{,}000 synthetic
  scenarios, 258 exhibited lower sum-check errors than the 1.2\% average
  error in the AR6 Scenarios Database.'' \emph{Rule:} the number of
  synthetic scenarios with electricity sum-check error $<1.2\%$ is
  $\ge 258$.
\item[\claim{X4} (VAE superiority).] ``VAE achieved higher category label
  transferring accuracy \dots compared to RCGAN and CGAN'' (Experiment~I,
  synthetic$\to$real, test population $=$ the 901-set). \emph{Rule:}
  accuracy(VAE) $>$ accuracy(RCGAN) and accuracy(VAE) $>$ accuracy(CGAN)
  in the same run.
\item[\claim{X5} (bias inheritance; secondary).] The paper positions
  synthetic generation as mitigating ``model-determined biases'' and
  ``non-random submission''. \emph{Rule:} train a source-model classifier
  (nine model groups, Dekker-style merge) on the real 901-set; apply to
  synthetic scenarios; X5 (inheritance) holds if the total-variation
  distance between the predicted model mix of the synthetic ensemble and
  the real ensemble's model shares is smaller than the distance to the
  uniform mix, per category. A baseline model-identifiability check on
  real data (held-out accuracy above the majority-class rate) gates the
  test.
\end{description}

\section{Factor space}\label{sec:factors}
Factors, paper values, and perturbation ranges are frozen in
\texttt{config/paper\_values.yaml} (archived with this deposit); the table
below summarizes. ``Expensive'' factors require retraining; ``cheap''
factors vary per trained model.

\begin{center}
\small
\begin{tabularx}{\textwidth}{@{}lXXl@{}}
\toprule
Factor & Paper value (provenance) & Range/levels & Cost\\
\midrule
G1 population & 1160-rule (D2; undeclared) & \{1202; 1160; 901; filter variants\} & exp.\\
\addlinespace
G2 aggregation & C1234/C56/C78 (Methods) & \{3-way; 4-way; 8-way; k-means labels\} & exp.\\
\addlinespace
G3 variable sel. & 5 features, RF+SHAP (SI S2) & completeness thr.\ \todo{0.6--0.9}; $n$ feat.\ \todo{4--8} & exp.\\
\addlinespace
G4 preprocessing & mean-imputation yr$\times$cat.; (0,1) scaling & imputation variants; scaling variants & exp.\\
\addlinespace
G5 architecture & VAE latent dim 10; ED Table 3 & latent \todo{4--20}; epochs $\times$\todo{0.5--2} & exp.\\
\addlinespace
G6 seed & seed 3; training non-reproducible (D1) & noise-floor replicates & exp.\\
\addlinespace
G7 sample size & 10{,}000/category & \todo{1{,}000--30{,}000} & cheap\\
\addlinespace
G8 transfer classifier & RF (code) & \{RF paper; RF alt; logistic; XGBoost\} & cheap\\
\addlinespace
G9 validation metrics & Fig.\ 3--5 statistics & OT spec (\S\ref{sec:estimators}); thresholds & cheap\\
\bottomrule
\end{tabularx}
\end{center}

\noindent The D4 documentation conflict (CGAN/RCGAN epochs: SI vs ED
Table~3) is handled as an explicit two-level factor nested in G5 unless
adjudicated from the shipped weights before freeze.

\section{Design and compute budget}\label{sec:design}
\begin{enumerate}[nosep]
\item \textbf{Noise floor.} $R_{0}=$ \todo{propose 20} seed-only
  replications at paper values for all factors. The floor $\delta_{0}$ is
  the \todo{95th percentile} of pairwise OT distances between replicate
  ensembles, per category.
\item \textbf{Outer design.} Sobol' quasi-random sample of size
  $N_{\mathrm{outer}}=$ \todo{propose 512} over the expensive factors
  (G1--G6), one training per point, environment pinned to A2's
  configuration.
\item \textbf{Inner design.} Per trained model, a full or fractional grid
  over cheap factors (G7--G9), $N_{\mathrm{inner}}=$ \todo{propose 8}.
\item \textbf{Stopping rule.} The budget is
  $N_{\mathrm{outer}}\times$(training cost); if bootstrap CIs on the
  leading OT indices have half-width $<$ \todo{0.05} at
  $N_{\mathrm{outer}}/2$, the design may stop early; it may not be
  extended beyond $N_{\mathrm{outer}}$ without a documented deviation.
\item \textbf{Outputs stored per run:} the generated ensemble (or
  sufficient statistics + per-marginal empirical quantiles), all claim
  evaluations X1a--X5, wall-clock, and the full factor vector.
\end{enumerate}

\section{Estimators}\label{sec:estimators}
\begin{itemize}[nosep]
\item \textbf{Ensemble-level sensitivity:} given-data optimal-transport
  sensitivity indices (Borgonovo--Chiani--Plischke class) on the synthetic
  ensemble as multivariate output. Distance: exact 1-D Wasserstein on
  each variable--year marginal, aggregated by \todo{mean}, plus sliced
  Wasserstein ($L=$ \todo{50} projections) on the joint; Sinkhorn
  ($\varepsilon=$ \todo{}) as robustness check. Partitioning per the
  given-data estimator with \todo{$M$} bins; discrete factors use exact
  level partitions (partition-collapse guard as in our prior OT-GSA
  implementation).
\item \textbf{Claim-level sensitivity:} for each rule X1a--X5, fraction of
  runs in which the claim holds, plus Sobol'-style total-order indices on
  the binary outcome, mirroring the companion analysis of Li P.-H.\ et
  al.\ (2023) for cross-paper comparability.
\item \textbf{Inference:} bootstrap with $R=1000$ resamples; 95\% CIs on
  all reported indices.
\end{itemize}

\section{Decision criteria}\label{sec:criteria}
\begin{itemize}[nosep]
\item A factor is \emph{influential} for the ensemble output if the lower
  95\% CI bound of its total OT index exceeds the noise floor
  $\delta_{0}$-equivalent index.
\item A claim is \emph{robust} if it holds in $\ge$ \todo{propose 80\%}
  of runs across the design; \emph{fragile} if in $<50\%$.
\item The study's headline question --- are the paper's findings
  properties of the framework or of the analyst's choices? --- is answered
  by the joint pattern of (i) claim fragility and (ii) the identity of the
  dominant factors (declared vs undeclared choices).
\end{itemize}

\section{Deviations policy}\label{sec:deviations}
Any departure from this protocol is logged in \texttt{DEVIATIONS.md} with
date, reason, and the analysis state at the time of the change. Results
affected by deviations are reported separately from the pre-registered
analysis. Adjudications of \todo{} placeholders completed \emph{before}
the Zenodo freeze are part of this protocol. The deposited PDF contains no
placeholders.

\section{Archive contents at deposit}
Protocol PDF, \texttt{DATA.md}, \texttt{DISCREPANCIES.md},
\texttt{config/paper\_values.yaml} (frozen), Phase-0 reports
(\texttt{results/*.json}, including the Phase~0B regeneration report),
v1.0.0/v1.0.1 diff, commit hash of the analysis harness. GSA results:
none, no claim rule evaluated at any factor setting (the paper's own
values included), no retraining performed, no noise-floor replicate run.

\end{document}
